\subsection{Client Entry Point}

\subsubsection{Purpose}
The purpose of the client entry point is to require every client module in order, and then initialize each client module in the same order.

\subsubsection{Functions}

% Optional overview paragraph
There are no functions

\subsubsection{Diagram}




\usetikzlibrary{positioning}

\begin{tikzpicture}

% Define your module names here
\def\modules{
    Core,
    Network,
    Database,
    Auth,
    Logging,
    UI,
    API,
    Cache,
    Config,
    Metrics,
    Storage,
    Security,
    Scheduler,
    Parser,
    Renderer
}

% Nodes
\node (get) [process, text width=6cm, align=left] {
Get all modules:\\
\small
\foreach [count=\i] \m in \modules {
    \i. \m\\
}
:(expandable)
};

\node (require) [process, below=of get.south, yshift=-5mm, text width=5cm, align=center] {
Require each module
};

\node (init) [process, below=of require.south, yshift=-5mm, text width=5cm, align=center] {
Initialize modules
};

% Arrows
\draw [arrow] (get) -- (require);
\draw [arrow] (require) -- (init);

\end{tikzpicture}

\subsubsection{Dependencies}

\subsubsection{Usage-example}

\subsubsection{Known issues and bottle-necks}

\subsubsection{Performance}